\section{Problem-Level Verification of LLM-Written QUBOs}
\label{sec:verification}

\subsection{The bitstring protocol and its failure modes}
QuantumQUBO Agent~\citep{mondal2026quantumqubo} verifies a generated QUBO on small test cases: an LLM converts each
human-written test case of a benchmark (instance, optimal value, an optimal solution in words) into a JSON case
with a variable count $n$ and the optimal \emph{bitstrings} in the formulation's own variable layout; code then
enumerates all $2^n$ bitstrings and accepts the QUBO if its minimizers overlap with the reference bitstrings.
The LLM that writes the reference bitstrings must therefore count the variables of the formulation, apply its
index map, and fill in auxiliary and slack variables correctly, all without executing anything.

We audited all 33 failures of the pipeline in a run with a local model (67/100 successes) by re-running the final
code of each failed benchmark on its test cases and inspecting the cases. Twelve failures were shape mismatches
between the generated $Q$ and the case's $n$. Comparing each case with the formulation's own variable map, the
\emph{test case} was wrong in at least 7 of the 12 (e.g.\ $n=9$ written for a model whose map gives
$4+5+5=14$ variables, while the code correctly built 14), the code was wrong in 2, and 3 were undetermined. Six
further failures were format mismatches between the case's instance and what the code read, and eleven were
argmin mismatches. One benchmark, quadratic assignment, is internally inconsistent: its stated optimal assignment
costs 20 while its stated optimum is 12 (the optimum of 12 is attained by a different assignment), so a correct
model fails whenever the reference bitstring follows the stated assignment.

\subsection{Problem-level verification}
We verify what a QUBO's ground states \emph{mean} instead of how they are laid out.
\begin{itemize}
  \item A test case holds the instance, the benchmark's stated optimal value (human-written) and, when given, a
        stated optimal solution in problem terms, plus one shared solution format; no variable count and no
        bitstrings.
  \item The Coder writes \texttt{decode(instance, x)} next to the QUBO builder.
  \item A Checker, written from the problem statement alone (it sees neither the formulation nor the code),
        writes \texttt{evaluate(instance, solution)}, which returns feasibility and the objective value.
  \item The Checker is trusted on a case only if it reproduces the stated optimal value on the stated optimal
        solution; untrusted cases are reported as \emph{unverifiable}, never as failures of the QUBO.
  \item A QUBO passes a case only if \emph{every} ground state decodes to a feasible solution with the stated
        optimal value; a ground state that is wrong could be returned by a solver.
  \item A case whose QUBO is too large to enumerate (here, more than 26 variables) is reported as unverifiable
        instead of being sent back to the Coder, which cannot change the test case.
\end{itemize}
The protocol removes three sources of error of the bitstring protocol: miscounted variables, wrongly filled
auxiliary variables, and copied inconsistencies in the benchmark text (the checker disagrees with the benchmark,
and the case becomes unverifiable). It also changes the acceptance rule in one direction on purpose: a QUBO whose
ground states include a non-optimal solution passed before (any overlap sufficed) and fails now.

\subsection{Results}
\Cref{tab:verification} compares the protocols on the 100 benchmarks of QuantumQUBO Agent with the same local model.
Five benchmarks failed under the bitstring protocol but succeeded in all three problem-level runs; three passed
under the bitstring protocol but failed all three problem-level runs, because their QUBOs have degenerate ground
states that include non-optimal solutions. Of the 16 unverifiable outcomes, 10 came from encodings too large to
enumerate (28 to 120 variables for tiny instances) and 6 from checkers that disagreed with the benchmark,
including the inconsistent quadratic-assignment case in all three runs.

\begin{table}[t]
\centering
\caption{Verification protocols on the 100 QuantumQUBO Agent benchmarks (seed 42). The authors' run used
remote models; ours use Qwen3.6-35B-A3B locally. Problem-level: three repeated runs.}
\label{tab:verification}
\small
\begin{tabular}{@{}lccc@{}}
\toprule
 & Successes & Unverifiable & Success among verifiable \\
\midrule
Bitstring, original authors' results file & 68 & -- & -- \\
Bitstring, local model & 67 & -- & -- \\
Problem-level, local model (3 runs) & 69 / 72 / 70 (mean 70.3) & 4 / 6 / 6 & 71.9 / 76.6 / 74.5\% \\
\bottomrule
\end{tabular}
\end{table}

\paragraph{Single runs are not reliable.}
All agents run at temperature zero, so repeated runs differ only through non-determinism in batched inference.
Even so, 27 of the 100 benchmarks changed outcome across the three identical problem-level runs (56 succeeded in
all three, 17 in none). Between the authors' results and our bitstring run, 25 benchmarks flipped while the totals
differed by one. Aggregate success rates are stable (standard deviation 1.5 points), but per-benchmark conclusions
from a single run are not, and comparisons between systems should use repeated runs.
